\documentclass[10pt,a4paper]{amsart}
\usepackage[margin=19mm]{geometry}
\usepackage{amsmath,amssymb,mathtools}
\usepackage[scr=rsfs,cal=euler]{mathalfa}
\usepackage{xfrac}
\usepackage[unicode,hidelinks,bookmarks=false]{hyperref}
\newcommand{\ie}{i.e.\ }
\newcommand{\cf}{cf.\ }
\newcommand{\resp}{respectively\ }
\newcommand{\Subsection}{\subsection}
\providecommand{\nobreakdash}{\nobreak\hbox{-}}
\setlength{\emergencystretch}{2em}
\multlinegap0pt
\usepackage{bm}
\usepackage{bbm}
\usepackage{dsfont}
\usepackage{xspace}
\usepackage{cancel}
\usepackage{mathtools}
\usepackage{amsthm}
\makeatletter
\@ifundefined{theorem}{\newtheorem{theorem}{Theorem}}{}
\@ifundefined{corollary}{\newtheorem{corollary}[theorem]{Corollary}}{}
\@ifundefined{proposition}{\newtheorem{proposition}[theorem]{Proposition}}{}
\makeatother
\DeclareMathOperator{\BGG}{BGG}
\DeclareMathOperator{\Kl}{Kl}
\DeclareMathOperator{\coh}{coh}
\DeclareMathOperator{\fp}{fp}
\DeclareMathOperator{\gen}{gen}
\DeclareMathOperator{\rigfp}{rig-fp}
\newcommand{\cH}{\mathcal{H}}
\newcommand{\cN}{\mathcal{N}}
\newcommand{\cQ}{\mathcal{Q}}
\newcommand{\cV}{\mathcal{V}}
\newcommand{\cX}{\mathcal{X}}
\newcommand{\dR}{\mathrm{dR}}
\newcommand{\rig}{\mathrm{rig}}
\newcommand{\wH}{\widetilde{H}}
\renewcommand{\leq}{\leqslant}
\title[On the Bloch-Kato conjecture for GSp(4)]{On the Bloch-Kato conjecture for GSp(4)}
\author{David Loeffler, Sarah Livia Zerbes}
\date{Source: arXiv:2003.05960v4 (CC BY 4.0)}
\begin{document}
\maketitle

\noindent\emph{Source and licence.} On the Bloch-Kato conjecture for GSp(4). Version arXiv:2003.05960v4 (CC BY 4.0), distributed under the Creative Commons Attribution 4.0 International licence, as recorded in the arXiv licence metadata for this exact version and shown on the abstract page https://arxiv.org/abs/2003.05960v4. Full rights notice: licenses/arxiv-2003.05960-CC-BY-4.0.txt.\par\smallskip

\noindent\textit{Editorial scope.} Counted unit: the Proposition whose source label is \texttt{prop:goodxi} in the article (source0.tex lines 4181--4188, source label \texttt{prop:goodxi}), delivered together with the article's own complete original proof of it (source0.tex lines 4189--4210, one \texttt{proof} environment of 22 lines). No mathematics is written, repaired or re-derived: statement and proof are the article's own text line for line. Required context retained uncounted: prop:kerZ (lines 3098--3103); eq:cohrigiso2 (lines 3207--3212); cor:fppairmapsS (lines 3520--3528). Every definition, display and label of those lines is retained unchanged. Omitted from this package and not redistributed: 1--3097, 3104--3206, 3213--3519, 3529--4180, 4211--6652. Nothing inside a retained excerpt is omitted or altered: every retained body is delivered exactly as the article's lines are written, so no omission receipt of any kind accompanies this package. Citations: the retained excerpts carry no citation command at all, so no bibliography entry of the article is transcribed and no reference is redistributed. Reference adaptation: none was needed and none was made; every reference inside the delivered unit resolves inside it, so no unresolved reference or citation is produced. Printed slips kept literal: no printed word, symbol, hypothesis or number of the counted statement or of its proof was altered. Layout and class: the article's own class is \texttt{ipart} with packages including amsmath, amscd, amssymb, mathrsfs, cancel, mathtools, fullpage, enumerate, thmtools, tabu, mdframed, fancyhdr, wrapfig, xspace; the wrapper is the lane's pinned \texttt{amsart} editorial wrapper at 10pt on A4, which restates the theorem-like environments the retained text needs and adds the article's own macro definitions that the retained text needs (\texttt{\textbackslash BGG}, \texttt{\textbackslash Kl}, \texttt{\textbackslash cH}, \texttt{\textbackslash cN}, \texttt{\textbackslash cQ}, \texttt{\textbackslash cV}, \texttt{\textbackslash cX}, \texttt{\textbackslash coh}, \texttt{\textbackslash dR}, \texttt{\textbackslash fp}, \texttt{\textbackslash gen}, \texttt{\textbackslash leq}, \texttt{\textbackslash rig}, \texttt{\textbackslash rigfp}, \texttt{\textbackslash wH}), with extra wrapper packages (bm, bbm, dsfont, xspace, cancel, mathtools). The delivered numbering is the unit's own. Assets: no retained span contains a figure environment, a graphics-inclusion command, a PDF-page inclusion, a table or a TikZ picture; no image file is copied into this package, the package declares no asset dependency and no image file is redistributed with it. Licence: version arXiv:2003.05960v4 of this article is distributed under the Creative Commons Attribution 4.0 International licence, as recorded in the arXiv licence metadata for this exact version and shown on the abstract page https://arxiv.org/abs/2003.05960v4. A full rights notice is delivered at licenses/arxiv-2003.05960-CC-BY-4.0.txt. Characters: every retained excerpt is pure ASCII, so the delivered engine, XeLaTeX with its default Latin Modern text font, reports no missing character.
  \begin{proposition}\label{prop:kerZ} As above, let $\cQ(T)=(1 - \tfrac{T}{\alpha})( 1 - \tfrac{T}{\beta})$. Then the class $\cQ(\Phi)\cdot \eta^{(2, m)}_{\coh,-D}$ lies in the kernel of $Z'$.
%   \begin{enumerate}
%   \item .
%   \item The class $\cQ(\Phi) \cdot \eta^{(2, m)}_{\coh, -D}$ lies in the kernel of $(Z')^3$.
%   \end{enumerate}
  \end{proposition}

    \begin{align}
    \alpha_{G, \rig, c}: \wH^{3}_{\dR,c}(\cX_{G,\Kl}^m,\cV,1+q)
    &\xrightarrow{\,\cong\,} \cH^{1,2}_c(\cX_{G,\Kl}^m,\BGG(\cV), 1+q) = H_c^2\left(\cX_{G,\Kl}^m,\cN^1\right)^{\nabla=0},\label{eq:cohrigiso1}\\
    \alpha_{G, \rig, c0}: \wH^{3}_{\dR,c0}(\cX_{G,\Kl}^{2,m},\cV,1+q)
    &\xrightarrow{\,\cong\,} \cH^{1,2}_{c0}(\cX_{G,\Kl}^{2,m},\BGG(\cV), 1+q) = H_{c0}^2\left(\cX_{G,\Kl}^{2,m},\cN^1\right)^{\nabla=0}.\label{eq:cohrigiso2}
    \end{align}

   \begin{corollary}
    \label{cor:fppairmapsS}
    If $0\leq q\leq r_2$, then the spectral sequence gives rise to an isomorphism
    \begin{align*}
     \alpha_{G, \rigfp, c0}: \mathscr{Z}^{1, 2}_{\fp,c0}(\cX_{G,\Kl}^{2,m}, \BGG(\cV),1+q; P)
     & \xrightarrow{\,\cong\,} \cH^{1, 2}_{\fp,c0}(\cX_{G,\Kl}^{2,m}, \BGG(\cV),1+q;P)\\
     & \xrightarrow{\,\cong\,} \wH^{3}_{\rigfp, c0}(\cX_{G,\Kl}^{2,m}, \cV,1+q; P).
    \end{align*}
   \end{corollary}

  \begin{proposition}\label{prop:goodxi}
 There exists $\xi\in  H^2_{c0}(\cX^{2,m}_{G,\Kl},\cN^0)$ with the following properties:
 \begin{enumerate}
  \item $\nabla\, \xi=\cQ_{1+q}(\Phi_{1+q})\eta_{\coh}^{(2,m)}$;
  \item $(U_2'-\lambda)\xi$ lies in the $U_2$-generalized eigen-subspace of $H^2_{\dR,c0}(\cX^{2,m}_{G,\Kl},\cV)$ with generalized eigenvalue $\lambda$;
  \item we have $Z'\cdot\xi=0$.
 \end{enumerate}
\end{proposition}

\begin{proof}

 \textbf{Step 1.}  We first show that there exists some $\xi$ such that $\nabla\,\xi=\cQ_{1+q}(\Phi_{1+q})\cdot \eta_{\coh}^{(2,m)}$.  By Proposition \ref{prop:kerZ}, we know that $Z'\circ \cQ(\Phi)\cdot \eta_{\coh,-D}^{(2,m)}=0$, which implies that
 \[ Z'\circ \cQ(\Phi)\cdot \eta_{\coh}^{2,m}=0. \]
 Recall that $\tilde\eta^{(2,m)}_{\rig}\in \wH^3_{\dR,c0}(\cX^{2,m}_{G,\Kl},\cV,1+q)$ is the preimage of $\eta_{\coh}^{(2,m)}$ under the isomorphism \eqref{eq:cohrigiso2}, so we deduce that
 \[ Z'\circ \cQ_{1+q}(\Phi_{1+q})\circ\iota ( \tilde\eta^{(2,m)}_{\rig})=0.\]
 Now recall that both $Z'$ and $\Phi$ commute with $U_2'$. Hence $\cQ_{1+q}(\Phi_{q})\circ\iota(\tilde\eta^{(2,m)}_{\rig})$ lies in the $(U_2'=\lambda)$ eigenspace of $ \wH^3_{\dR,c0}(\cX^{2,m}_{G,\Kl},\cV)$, and since $Z'\circ\Phi=p^{r_2+1}\, U_2'$, the restriction of $Z'$ to this eigenspace is a bijection. We deduce that
 \[  \cQ_{1+q}(\Phi_{1+q}) \circ \iota (\tilde\eta^{(2,m)}_{\rig})=0.\]
 We can hence lift $\tilde\eta^{(2,m)}_{\rig}$ to a class
 \[ \tilde\eta^{(2,m)}_{\rigfp} \in \wH^3_{\rigfp,c0}(\cX^{2,m}_{G,\Kl},\cV,1+q;\cQ_{1+q})   \]
 which lies in the $\Pi'_f$-eigenspace for the spherical Hecke operators. By Corollary \ref{cor:fppairmapsS}, this class corresponds to a coherent fp-pair, which has the required form.
 \vspace{1ex}

 \textbf{Step 2.} Note that
 \begin{itemize}
  \item we have  $(U_2'-\lambda)\xi\in H^2_{c0}(\cX^{2,m}_{G,\Kl},\cN^0)^{\nabla=0}$, since $U_2'\,\eta_{\coh}^{(2,m)}=\lambda\eta_{\coh}^{(2,m)}$,;
  \item we have $Z'\cdot\xi\in H^2_{c0}(\cX^{2,m}_{G,\Kl},\cN^0)^{\nabla=0}$, by Proposition \ref{prop:kerZ} (1).
 \end{itemize}
 Now $H^2_{c0}(\cX^{2,m}_{G,\Kl},\cN^0)^{\nabla=0}\cong H^2_{\dR,c0}(\cX^{2,m}_{G,\Kl},\cV)$ is finite-dimensional, so by applying a suitable projector we can assume without loss of generality that both $(U_2'-\lambda)\xi$ amd $Z'\cdot \xi$ lie in the $U_2$-generalized eigen-subspace of $H^2_{\dR,c0}(\cX^{2,m}_{G,\Kl},\cV)$ with generalized eigenvalue $\lambda$ (we use here that $U_2'$ commutes with $Z'$); denote this subspace by
 \[ H^2_{\dR,c0}(\cX^{2,m}_{G,\Kl},\cV)[U_2'=\lambda]^{\gen}.\]
 Now since $Z'\circ\Phi=p^{r_2+1}\, U_2'$, the restriction of $Z'$ to $H^2_{\dR,c0}(\cX^{2,m}_{G,\Kl},\cV)[U_2'=\lambda]^{\gen}$ is a bijection, so there exists $\nu\in H^2_{\dR,c0}(\cX^{2,m}_{G,\Kl},\cV)[U_2'=\lambda]^{\gen}$ such that $Z'\cdot \nu=Z'\cdot \xi$. Replacing $\xi$ by $\xi-\nu$ proof the result.
\end{proof}

\end{document}
